% !TeX program = xelatex
%% =====================================================================
%%  ใบกิจกรรมในชั้นเรียน (Worksheet) บทที่ 1 : ตรรกศาสตร์ --- สัปดาห์ที่ 3
%%  เนื้อหา: ตัวบ่งปริมาณ (Quantifiers) — สอดคล้องกับสไลด์
%%  ระดับง่าย-กลาง (6 ข้อ, 15-20 นาที) — มีช่องว่างให้เขียน
%% =====================================================================
\documentclass[a4paper,12pt]{article}
\usepackage{fontspec}
\usepackage{amsmath,amssymb}
\usepackage[margin=2cm]{geometry}
\usepackage{enumitem}
\usepackage{array}
\usepackage{xcolor}
\usepackage{colortbl}

\XeTeXlinebreaklocale "th_TH"
\XeTeXlinebreakskip = 0pt plus 1pt
\setmainfont[Script=Thai,Scale=1.25]{TH Sarabun New}
\definecolor{navy}{HTML}{1F3A93}
\definecolor{Tgreen}{HTML}{2E8B57}
\definecolor{Fred}{HTML}{C0392B}
\definecolor{gold}{HTML}{B7950B}
\definecolor{lightnavy}{HTML}{EAEEF7}
\renewcommand{\arraystretch}{1.4}
\setlist{topsep=4pt,itemsep=6pt}

%% เส้นว่างให้เขียนคำตอบ (#1 = จำนวนบรรทัด)
\newcounter{bl}
\newcommand{\blank}[1]{\par\vspace{0.2em}%
  \setcounter{bl}{0}\loop\ifnum\value{bl}<#1\relax
  \noindent\textcolor{gray!60}{\hrulefill}\par\vspace{1.0em}\stepcounter{bl}\repeat}

\newcommand{\Tcell}{\textcolor{Tgreen}{\textbf{T}}}
\newcommand{\Fcell}{\textcolor{Fred}{\textbf{F}}}
\newcommand{\A}{\textcolor{navy}{$\forall$}}
\newcommand{\E}{\textcolor{gold}{$\exists$}}

\begin{document}
\thispagestyle{empty}

%% ---------- หัวกระดาษ ----------
\begin{center}
  {\Large\bfseries\color{navy} ใบกิจกรรมในชั้นเรียน บทที่ 1 : ตัวบ่งปริมาณ (สัปดาห์ที่ 3)}\\[0.2em]
  {\normalsize รายวิชา 4091606 คณิตศาสตร์สำหรับคอมพิวเตอร์}
\end{center}
\vspace{0.3em}
\noindent\begin{tabular}{@{}p{0.5\textwidth} p{0.22\textwidth} p{0.22\textwidth}@{}}
  ชื่อ-สกุล: \dotfill & รหัส: \dotfill & กลุ่ม: \dotfill
\end{tabular}
\hrule
\vspace{0.5em}
\noindent\textbf{คำชี้แจง:} ให้นักศึกษาทำกิจกรรม 6 ข้อต่อไปนี้ลงในใบกิจกรรม \textbf{ภายใน 15-20 นาที}
แสดงวิธีคิดให้ชัดเจน แล้วร่วมเฉลยและอภิปรายพร้อมกันในชั้นเรียน
\vspace{0.5em}

%% ---------- ข้อ 1 : ประพจน์เปิด ----------
\noindent\textbf{ข้อ 1: ประพจน์เปิด (Open Proposition)} \hfill \textit{(2 นาที)} \\
กำหนดประพจน์เปิด $P(x)$: ``$x$ เป็นเลขคู่'' จงหาค่าความจริงเมื่อแทนค่า $x$ ต่อไปนี้
\begin{center}
\begin{tabular}{|c|c|c|c|c|}
  \hline
  \textbf{แทนค่า} & $x = 2$ & $x = 3$ & $x = 8$ & $x = 5$ \\ \hline
  \textbf{$P(x)$ จริง/เท็จ} & & & & \\ \hline
\end{tabular}
\end{center}

%% ---------- ข้อ 2 : แปลงเป็นสัญลักษณ์ ----------
\noindent\textbf{ข้อ 2: แปลงประโยคเป็นสัญลักษณ์ตัวบ่งปริมาณ} \hfill \textit{(3 นาที)} \\
จงเขียนประโยคต่อไปนี้ในรูป \A$x\,[P(x)]$ หรือ \E$x\,[P(x)]$
\quad(เอกภพ = \textbf{เซตของนักศึกษา})
\begin{enumerate}[label=\arabic*)]
  \item ``นักศึกษา\textbf{ทุกคน}ส่งการบ้าน''\\
        \textbf{คำตอบ:} \blank{1}
  \item ``\textbf{มี}นักศึกษา\textbf{บางคน}ได้เกรด A''\\
        \textbf{คำตอบ:} \blank{1}
  \item ``\textbf{ไม่มี}นักศึกษาคนใดขาดสอบ''\\
        \textbf{คำตอบ:} \blank{1}
\end{enumerate}

%% ---------- ข้อ 3 : หาค่าความจริง ----------
\noindent\textbf{ข้อ 3: หาค่าความจริงบนเอกภพจำกัด} \hfill \textit{(3 นาที)} \\
กำหนดเอกภพ $U = \{1, 2, 3, 4\}$ จงหาค่าความจริง พร้อมแสดงวิธีตรวจสอบ
\begin{enumerate}[label=\arabic*)]
  \item \A$x \in U\,[x + 1 \le 5]$ \hfill
        \textbf{ค่าความจริง:} \dotfill\\
        \textbf{วิธีตรวจ:} \blank{1}
  \item \E$x \in U\,[x^2 = 9]$ \hfill
        \textbf{ค่าความจริง:} \dotfill\\
        \textbf{วิธีตรวจ:} \blank{1}
\end{enumerate}

%% ---------- ข้อ 4 : ตัวอย่างค้าน ----------
\noindent\textbf{ข้อ 4: ยกตัวอย่างค้าน (Counterexample)} \hfill \textit{(3 นาที)} \\
กำหนดเอกภพ $U = \{1, 2, 3, 4, 5, 6, 7, 8, 9, 10\}$
\begin{enumerate}[label=\arabic*)]
  \item \A$x\,[x > 5]$ เป็นจริงหรือเท็จ? ถ้าเท็จ จงยก \textbf{ตัวอย่างค้าน} 1 ตัว\\
        \textbf{ค่าความจริง:} \dotfill \quad \textbf{ตัวอย่างค้าน:} \blank{1}
  \item \E$x\,[x \text{ เป็นจำนวนเฉพาะ}]$ เป็นจริงหรือเท็จ? ถ้าจริง จงยก \textbf{ตัวอย่างยืนยัน} 1 ตัว\\
        \textbf{ค่าความจริง:} \dotfill \quad \textbf{ตัวอย่างยืนยัน:} \blank{1}
\end{enumerate}

%% ---------- ข้อ 5 : นิเสธ ----------
\noindent\textbf{ข้อ 5: เขียนนิเสธของตัวบ่งปริมาณ} \hfill \textit{(4 นาที)} \\
จงเขียนนิเสธ ทั้ง\textbf{รูปสัญลักษณ์}และ\textbf{ภาษาไทย}
\quad\textit{(สูตร: $\neg$(\A$x\,[P(x)]$) $\equiv$ \E$x\,[\neg P(x)]$)}
\begin{enumerate}[label=\arabic*)]
  \item ประโยค: ``นักศึกษาทุกคนมาเรียนตรงเวลา''\\
        \textbf{สัญลักษณ์เดิม:} \dotfill\\
        \textbf{นิเสธ (สัญลักษณ์):} \dotfill \quad \textbf{นิเสธ (ไทย):} \blank{1}
  \item ประโยค: ``มีคอมพิวเตอร์บางเครื่องที่เสีย''\\
        \textbf{สัญลักษณ์เดิม:} \dotfill\\
        \textbf{นิเสธ (สัญลักษณ์):} \dotfill \quad \textbf{นิเสธ (ไทย):} \blank{1}
\end{enumerate}

%% ---------- ข้อ 6 : ตัวบ่งปริมาณซ้อน ----------
\noindent\textbf{ข้อ 6: ตัวบ่งปริมาณซ้อน (Nested)} \hfill \textit{(5 นาที)} \\
กำหนด $F(x,y)$: ``$x$ เป็นเพื่อนกับ $y$'' จงอธิบายความหมายของประพจน์ต่อไปนี้เป็นภาษาไทย
\begin{enumerate}[label=\arabic*)]
  \item \A$x,\,$\E$y\,[F(x,y)]$\\
        \textbf{ความหมาย:} \blank{1}
  \item \E$x,\,$\A$y\,[F(x,y)]$\\
        \textbf{ความหมาย:} \blank{1}
  \item ทั้งสองข้อข้างต้น \textbf{มีความหมายเหมือนกันหรือไม่}? เพราะเหตุใด\\
        \blank{1}
\end{enumerate}

\vspace{0.5em}
\begin{center}
  \color{navy}\textbf{--- จบใบกิจกรรม ---}
\end{center}

\end{document}
